% page 234 du fac-simile (image p-233 du scan)
% bloc 162.6 x 229.0 mm ; marge gauche 8.0 mm
\pgc{-12.700mm}{0.000mm}{
\l{\marge[84.642mm]{24.468mm}{\ornamento{6.961mm}{ornements/letroj/litero-I.png}}\cel{5}226}
\l{}
\l{}
\l{}
\l{}
\l{}
\l{}
\l{}
\l{}
\l{}
\l{\sou{-i-}. Sufixo qua indikas \textquotedbl{}lando, regiono, domeno qua dependas}
\l{\cel{3}de... e, konseque, la resortiso di, la autoritato di...\textquotedbl{}}
\l{}
\l{\sou{-i.}\cel{3}Dezinenco plurala di la nomi e di la pronomi.}
\l{}
\l{\sou{iambo}. I. (\sou{metriko antiqua)}.\cel{2}I. Pedo ek du silabi, la unesma,}
\l{\cel{3}kurta ; la duesma, longa. - 2. Verso satirala, qua konsistis}
\l{\cel{3}prime ek sis iambi. - II. Versajo satira. - DEFIRS.}
\l{}
\l{\sou{ibe.}\cel{4}En ta loko (quan onu indikas), kontraste kun olta en}
\l{\cel{3}qua onu esas.}
\l{}
\l{\sou{ibiso}.\cel{2}(\sou{zool.}) Ucelo steltiera, longa-beka, di qua un speco}
\l{\cel{3}esis sakra, che la Egiptiani antiqua. - L.\cel{1}\sou{ibis}. - DEFIRS.}
\l{}
\l{\sou{ica}.\cel{3}ca.}
\l{}
\l{\sou{-id-}.\cel{2}Sufixo qua signifikas \textquotedbl{}homo,decendanto di...\textquotedbl{}}
\l{}
\l{\sou{idanto}. (\sou{biol.}) Segun la teorio di Waismann : partikuli qui}
\l{\cel{3}reprezentas e transmisas atavisme la partikularaji di la}
\l{\cel{3}strukturo di la enti vivoza.}
\l{}
\l{\sou{ideo.}\cel{2}Imajo, en spirito, de ula ento o de ula eso-maniero.-}
\l{\cel{3}DEFIRS.}
\l{}
\l{\sou{identa}. Di qua la naturo esas absolute sama kam olta di altra}
\l{\cel{3}kozo. - DEFIRS.}
\l{}
\l{\sou{ideografio.}\cel{2}Reprezento di la idei per signi qui esas la imajo}
\l{\cel{3}de la objekto. - DEFIRS.}
\l{}
\l{\sou{ideologio}.\cel{2}(\sou{filoz.}) La cienco pri la origino, la formac(es)o}
\l{\cel{3}di la idei. - DEFIRS.}
\l{}
\l{\sou{idilio.}\cel{2}I. Mikra poemo pastorala. - II. Amoro-naraco pastorala.}
\l{\cel{3}- DEFIRS.}
\l{}
\l{\sou{idioblasto.}\cel{2}(\sou{bot.)}\cel{2}(\sou{segun Sachs}), celuli neregulale branchi-}
\l{\cel{3}gita, izolita meze di tisuo tre diferanta,}
\l{}
\l{\sou{idiomo}. I. Linguo di naciono. - II. Dialekto di provinco. -}
\l{\cel{3}DEFIRS.}
}
